\documentclass[conference]{IEEEtran}
\usepackage{cite}
\usepackage{amsmath,amssymb,amsfonts}
\usepackage{algorithmic}
\usepackage{graphicx}
\usepackage{textcomp}
\usepackage{xcolor}
\def\BibTeX{{\rm B\kern-.05em{\sc i\kern-.025em b}\kern-.08em
    T\kern-.1667em\lower.7ex\hbox{E}\kern-.125emX}}
\begin{document}

\title{COS 711 Assignment 2\\
{\footnotesize Supervised Neural Network Training}
}

\author{\IEEEauthorblockN{Jeffrey Russell (u16010648)}
\IEEEauthorblockA{\textit{EBIT} \\
\textit{University of Pretoria}\\
Pretoria, South Africa \\
u16010648@tuks.co.za}
}

\maketitle

\begin{abstract}
In this project we were instructed to implement a neural network with
automatic architecture selection and curriculum learning on two wine quality
datasets. There were numerous hurdles to overcome such as the nature of the data
and the imbalance of the classes given. The final results were promising with 
an accuracy of approximately 70\% finally being achieved on the red wine and
<insert here> on the white wine.
\end{abstract}

\begin{IEEEkeywords}
neural networks, curriculum learning
\end{IEEEkeywords}

\section{Introduction}
In this project we were given a dataset containing a number of wine traits
(ph, density, etc.) and the subjective quality assigned by a human. My goal was
to create a neural network with certain added extras that would be able to predict
the quality of the wine given the set of wine traits of any particular wine. The
majority of the work fell into four different pieces of this program: preprocessing,
manual parameter tuning, auto architecture selection, and a curriculum learning algorithm.

The results were promising, achieving good accuracy on the two datasets, but
they did seem a bit underwhelming considering the difficulty to produce them.

\section{Background}
The preprocessing of the data was fairly simple because Python has so many built-in
packages for doing data manipulation and processing. The pre processing involved some
statistical analysis and the PCA algorithm to remove useless properties of the wine, 
and reduce the dimensionality of the problem for the NN.
\\
After the preprocessing, the next step was to create the neural network itself.
I used Tensorflow for the neural network itself with categorical cross entropy
as the loss function (this isn't the full story on the loss function, more on that later).
The preprocessing step took care of the test/generalization/global split. The 
parameter tuning that I did to come up with an optimal neural network architecture
and training strategy was fairly emperical, trying different values and seeing
what worked and what didn't.
\\
Insert auto architecture selection piece here
\\
The curriculum learning part of the project ended up being fairly simple, though
it was very long running. I created a GA which decides on the order
of the samples shown to the neural network in a number of batches.
It progressively shows the NN more and more samples until in the final batch
all of the samples are present.


\section{Experimental set-up}
I've split the experimental set-up into the four different parts of the project.

\subsection{Preprocessing}
I loaded the data through pandas from the csv files and then used scipy and numpy to split up the data and do the preprocessing.
The data was first sent through a z-score estimation function (provided by scipy.stats) which assumes a normal distribution of each input variable and then gives each value a z-score based on the mean and standard deviation within the input variable.
The assumption of a normal distribution seemed justifiable since since most of the qualities are based on biological processes within the fermentation of wine that should present normal distribution curves when examined across a sufficiently large sample size. A couple of the input variables were fairly skewed to one side or the other, but the removal of outliers and increase in variation that the z-score introduced would seem to outweigh the possible problems with this approach.
I then removed all rows in the data which had a column that was more than three standard deviations away from the mean. This dealt with outliers nicely and made the rest of the process much simpler. I also reduced the range of the values to between 0 and 1. In order to deal with class imbalance I decided to use a cost-sensitive loss function based on how many samples there were of that kind, the exact equation for the loss is described in the neural network architecture section.
The next step for the preprocessing was to apply a PCA to the dataset in order to throw away the least useful input variables. I decided to throw away two of the input variables since that seemed to work well. That reduced the dimensionality of the input from 11 down to 9.
The final step was to split the data into training/testing and global generalization. The percentages I decided on was 70\% for training and then 15\% each for testing and global generalization. This split was decided on because it seemed to give good results without cutting down too much on the testing size. This was all done with the red wine dataset in mind which is a bit smaller so that influenced the decision a bit. Because all the processing steps happened before the split between test and training data, there might have been some latent information based on the preprocessing that might have slightly increased the accuracy, but I don't think that it would be enough to warrant a different approach.

\subsection{Neural network}
The neural network was created in Tensorflow with Keras as the backend. I did a lot of experimentation trying to work out which particular parameters ended up creating the best and quickest convergence. Not everything that I did I have graphs for, but the most important stuff I do. The loss function that I used was categorical cross entropy (ish, I'll discuss the loss function below since it has some interesting quirks and rationale behind it). I decided on this because it seems to be the standard for measuring loss on categorical data across the sources that I reviewed.
The input layer uses the tanh activation function. The hidden layers ended up just being a set of five layers with 9 neurons each all with activation function tanh. The output layer has 10 neurons with sigmoid activation. The input activation function is tanh because as soon as I tried any other activation function I immediately got a NaN loss in tensorflow. It seems to be because of the range of the output combined with the loss function which is categorical crossentropy which has issues with predictions of exactly zero since it has to do the log of that which causes undefined math operations. The hidden layers are all tanh since that had the best accuracy at a low number of epochs and a comparable accuracy further on. See fig 1 and 2 which compared relu, tanh, sigmoid and linear over 500 epochs running each model 100 times. Interestingly sigmoid is absolutely terrible up until about generation 300, and I have no idea why this might be.
Because of the data shown in fig 1 and 2 I decided to have a minimum of 200 epochs with an ending function. This function takes a look at the last five epochs. For each epoch it considers the loss function result on the test data and compares that to the result for the current epoch. If 3 out of the most recent five have lower losses than the current epoch then the training is aborted. This exact function seems to strike an excellent balance between stopping too early based on a false positive and stopping too late because it didn't see the trend. As soon as the test loss starts increasing consistently it stops the training. This tends to happen right around the 200 epoch mark


The first batch shown to the neural network has a fifth of the samples and takes a fifth of the total number of epochs. The following batch adds the next fifth of samples given by the GA and shows those in addition to the first batch to the neural network. This continues until for the last fifth of the training the neural network is training on the whole set. This seems to work quite well, and it gives an easily obtainable difficulty for samples depending on where they tend to end up in the order that the GA produces. I stumbled upon this way of creating the curriculum after reading through a number of different methodologies and realizing that this particular method seemed to fit well with the type of data and the other pieces of the puzzle that make up the classifier. 

\subsection{Abbreviations and Acronyms}\label{AA}
Define abbreviations and acronyms the first time they are used in the text, 
even after they have been defined in the abstract. Abbreviations such as 
IEEE, SI, MKS, CGS, ac, dc, and rms do not have to be defined. Do not use 
abbreviations in the title or heads unless they are unavoidable.

\subsection{Units}
\begin{itemize}
\item Use either SI (MKS) or CGS as primary units. (SI units are encouraged.) English units may be used as secondary units (in parentheses). An exception would be the use of English units as identifiers in trade, such as ``3.5-inch disk drive''.
\item Avoid combining SI and CGS units, such as current in amperes and magnetic field in oersteds. This often leads to confusion because equations do not balance dimensionally. If you must use mixed units, clearly state the units for each quantity that you use in an equation.
\item Do not mix complete spellings and abbreviations of units: ``Wb/m\textsuperscript{2}'' or ``webers per square meter'', not ``webers/m\textsuperscript{2}''. Spell out units when they appear in text: ``. . . a few henries'', not ``. . . a few H''.
\item Use a zero before decimal points: ``0.25'', not ``.25''. Use ``cm\textsuperscript{3}'', not ``cc''.)
\end{itemize}

\subsection{Equations}
Number equations consecutively. To make your 
equations more compact, you may use the solidus (~/~), the exp function, or 
appropriate exponents. Italicize Roman symbols for quantities and variables, 
but not Greek symbols. Use a long dash rather than a hyphen for a minus 
sign. Punctuate equations with commas or periods when they are part of a 
sentence, as in:
\begin{equation}
a+b=\gamma\label{eq}
\end{equation}

Be sure that the 
symbols in your equation have been defined before or immediately following 
the equation. Use ``\eqref{eq}'', not ``Eq.~\eqref{eq}'' or ``equation \eqref{eq}'', except at 
the beginning of a sentence: ``Equation \eqref{eq} is . . .''

\subsection{\LaTeX-Specific Advice}

Please use ``soft'' (e.g., \verb|\eqref{Eq}|) cross references instead
of ``hard'' references (e.g., \verb|(1)|). That will make it possible
to combine sections, add equations, or change the order of figures or
citations without having to go through the file line by line.

Please don't use the \verb|{eqnarray}| equation environment. Use
\verb|{align}| or \verb|{IEEEeqnarray}| instead. The \verb|{eqnarray}|
environment leaves unsightly spaces around relation symbols.

Please note that the \verb|{subequations}| environment in {\LaTeX}
will increment the main equation counter even when there are no
equation numbers displayed. If you forget that, you might write an
article in which the equation numbers skip from (17) to (20), causing
the copy editors to wonder if you've discovered a new method of
counting.

{\BibTeX} does not work by magic. It doesn't get the bibliographic
data from thin air but from .bib files. If you use {\BibTeX} to produce a
bibliography you must send the .bib files. 

{\LaTeX} can't read your mind. If you assign the same label to a
subsubsection and a table, you might find that Table I has been cross
referenced as Table IV-B3. 

{\LaTeX} does not have precognitive abilities. If you put a
\verb|\label| command before the command that updates the counter it's
supposed to be using, the label will pick up the last counter to be
cross referenced instead. In particular, a \verb|\label| command
should not go before the caption of a figure or a table.

Do not use \verb|\nonumber| inside the \verb|{array}| environment. It
will not stop equation numbers inside \verb|{array}| (there won't be
any anyway) and it might stop a wanted equation number in the
surrounding equation.

\subsection{Some Common Mistakes}\label{SCM}
\begin{itemize}
\item The word ``data'' is plural, not singular.
\item The subscript for the permeability of vacuum $\mu_{0}$, and other common scientific constants, is zero with subscript formatting, not a lowercase letter ``o''.
\item In American English, commas, semicolons, periods, question and exclamation marks are located within quotation marks only when a complete thought or name is cited, such as a title or full quotation. When quotation marks are used, instead of a bold or italic typeface, to highlight a word or phrase, punctuation should appear outside of the quotation marks. A parenthetical phrase or statement at the end of a sentence is punctuated outside of the closing parenthesis (like this). (A parenthetical sentence is punctuated within the parentheses.)
\item A graph within a graph is an ``inset'', not an ``insert''. The word alternatively is preferred to the word ``alternately'' (unless you really mean something that alternates).
\item Do not use the word ``essentially'' to mean ``approximately'' or ``effectively''.
\item In your paper title, if the words ``that uses'' can accurately replace the word ``using'', capitalize the ``u''; if not, keep using lower-cased.
\item Be aware of the different meanings of the homophones ``affect'' and ``effect'', ``complement'' and ``compliment'', ``discreet'' and ``discrete'', ``principal'' and ``principle''.
\item Do not confuse ``imply'' and ``infer''.
\item The prefix ``non'' is not a word; it should be joined to the word it modifies, usually without a hyphen.
\item There is no period after the ``et'' in the Latin abbreviation ``et al.''.
\item The abbreviation ``i.e.'' means ``that is'', and the abbreviation ``e.g.'' means ``for example''.
\end{itemize}
% An excellent style manual for science writers is \cite{b7}.

\subsection{Authors and Affiliations}
\textbf{The class file is designed for, but not limited to, six authors.} A 
minimum of one author is required for all conference articles. Author names 
should be listed starting from left to right and then moving down to the 
next line. This is the author sequence that will be used in future citations 
and by indexing services. Names should not be listed in columns nor group by 
affiliation. Please keep your affiliations as succinct as possible (for 
example, do not differentiate among departments of the same organization).

\subsection{Identify the Headings}
Headings, or heads, are organizational devices that guide the reader through 
your paper. There are two types: component heads and text heads.

Component heads identify the different components of your paper and are not 
topically subordinate to each other. Examples include Acknowledgments and 
References and, for these, the correct style to use is ``Heading 5''. Use 
``figure caption'' for your Figure captions, and ``table head'' for your 
table title. Run-in heads, such as ``Abstract'', will require you to apply a 
style (in this case, italic) in addition to the style provided by the drop 
down menu to differentiate the head from the text.

Text heads organize the topics on a relational, hierarchical basis. For 
example, the paper title is the primary text head because all subsequent 
material relates and elaborates on this one topic. If there are two or more 
sub-topics, the next level head (uppercase Roman numerals) should be used 
and, conversely, if there are not at least two sub-topics, then no subheads 
should be introduced.

\subsection{Figures and Tables}
\paragraph{Positioning Figures and Tables} Place figures and tables at the top and 
bottom of columns. Avoid placing them in the middle of columns. Large 
figures and tables may span across both columns. Figure captions should be 
below the figures; table heads should appear above the tables. Insert 
figures and tables after they are cited in the text. Use the abbreviation 
``Fig.~\ref{fig}'', even at the beginning of a sentence.

\begin{table}[htbp]
\caption{Table Type Styles}
\begin{center}
\begin{tabular}{|c|c|c|c|}
\hline
\textbf{Table}&\multicolumn{3}{|c|}{\textbf{Table Column Head}} \\
\cline{2-4} 
\textbf{Head} & \textbf{\textit{Table column subhead}}& \textbf{\textit{Subhead}}& \textbf{\textit{Subhead}} \\
\hline
copy& More table copy$^{\mathrm{a}}$& &  \\
\hline
\multicolumn{4}{l}{$^{\mathrm{a}}$Sample of a Table footnote.}
\end{tabular}
\label{tab1}
\end{center}
\end{table}

\begin{figure}[htbp]
\centerline{\includegraphics{fig1.png}}
\caption{Example of a figure caption.}
\label{fig}
\end{figure}

Figure Labels: Use 8 point Times New Roman for Figure labels. Use words 
rather than symbols or abbreviations when writing Figure axis labels to 
avoid confusing the reader. As an example, write the quantity 
``Magnetization'', or ``Magnetization, M'', not just ``M''. If including 
units in the label, present them within parentheses. Do not label axes only 
with units. In the example, write ``Magnetization (A/m)'' or ``Magnetization 
\{A[m(1)]\}'', not just ``A/m''. Do not label axes with a ratio of 
quantities and units. For example, write ``Temperature (K)'', not 
``Temperature/K''.

\section*{Acknowledgment}

The preferred spelling of the word ``acknowledgment'' in America is without 
an ``e'' after the ``g''. Avoid the stilted expression ``one of us (R. B. 
G.) thanks $\ldots$''. Instead, try ``R. B. G. thanks$\ldots$''. Put sponsor 
acknowledgments in the unnumbered footnote on the first page.
\cite{CORTEZ2009547}


\bibliographystyle{IEEEtran}
\bibliography{IEEEabrv,u16010648assignment2}

\vspace{12pt}
\color{red}
IEEE conference templates contain guidance text for composing and formatting conference papers. Please ensure that all template text is removed from your conference paper prior to submission to the conference. Failure to remove the template text from your paper may result in your paper not being published.

\end{document}
